\section{Inherited diagonal-calculus calibration}\label{sec:diagonal}
This appendix reproduces the diagonal-interface argument from the supplied
diagonal-calculus report, as carried in the parent archive. Its historical
physical implementation paragraph is omitted; the195-state construction in
this report preserves the later separated-head implementation. Replace the equality-vacuity component of the gate by the diagonal clause
\begin{equation}\label{eq:gamma}
 \Gamma(B)=I(B,U(B,B)),\qquad\text{that is,}\qquad B\longrightarrow\forall x_B\,B.
\end{equation}
Here \(B\) is both a formula code and a variable index. This clause is sound because every variable index in the decoded formula \(B\) is strictly below \(B\). Code zero is truth and presents no exceptional free variable.

For the implementation included in this archive, replace the old gate
\[
 F(a,b,c)=\neg\bigl(\Delta_c(a,b)\to(V(a,b,c)\to Q(a,b,c))\bigr)
\]
by
\begin{equation}\label{eq:new-gate}
 F_\Gamma(a,b,c)=\neg\bigl(\Gamma(c)\to(V(a,b,c)\to Q(a,b,c))\bigr).
\end{equation}
Keep the old guard \(c\ne a,b\), the six other heads, the restricted propositional head, generalization, and modus ponens. Thus the membership-vacuity part of the gate is retained. Removing that part as well would invalidate the reverse-rewriting argument below.

\subsection{What can be extracted from the new gate}
Classical implicational reasoning gives both
\[
 F_\Gamma(a,b,c)\to\Gamma(c),\qquad
 F_\Gamma(a,b,c)\to G(a,b,c),
 \quad G=\neg(V\to Q).
\]
Every diagonal instance \(\Gamma(B)\) is derivable: take \(c=B\) and \(a=b=B+1\). Both guard comparisons succeed. The extracted old gate \(G\) continues to give unrestricted membership vacuity at labels different from the atom's indices, and the original pointed-separation schema. Coincident labels \(a=b\) are permitted by the guard and have already been accounted for in the arithmetic interpretation.

The retained propositional and quantifier heads give the same basic implicational reasoning, universal monotonicity, and the existential-elimination rule for a conclusion known to be vacuous at the quantified label. Membership vacuity extends to the quantifier-free membership matrices used below by the original negation and implication arguments. It also extends through a universal quantifier whose binder differs from the new vacuity label, by quantifier commutation. No equality-vacuity premise is used in those extensions.

\subsection{Recovering reverse rewriting}
The old proof used equality vacuity to discharge an existential witness from a conclusion containing an equality atom. The replacement is to make the separator property the negation of that entire conclusion.

Let \(a\ne b\). Let \(\chi(t)\) be one of
\(M(d,t)\), \(M(t,d)\), or \(M(t,t)\), with \(d\ne a\) in the first two cases. The direct membership-rewriting heads give
\[
 E(v,y)\to\bigl(\chi(v)\to\chi(y)\bigr).
\]
Put
\[
 C=E(a,b)\to\bigl(\chi(b)\to\chi(a)\bigr),
 \qquad \varphi=\neg C,\qquad c=\ulcorner\varphi\urcorner.
\]
The code \(c\) exceeds all variable labels in \(C\), so \(c\ne a,b\); in the first two families it also gives \(c\ne d\). There is no property-code collision requiring a separate padding branch.

Use pointed separation at \((a,a,c)\). It gives
\begin{equation}\label{eq:witness}
 \exists x_c\,(S\wedge P),\qquad
 P=\varphi\to M(c,a),\qquad
 S\to U\bigl(a,M(c,a)\to\varphi\bigr).
\end{equation}
Propositionally, \(\varphi\) implies \(E(a,b)\), \(\chi(b)\), and \(\neg\chi(a)\). In particular, weakening the last consequent in~\eqref{eq:witness} gives
\[
 S\to U\bigl(a,M(c,a)\to\neg\chi(a)\bigr).
\]
The original negative-property specialization lemma now applies. It uses the equality witness head, the direct membership rewrites, and membership-only vacuity, and yields
\[
 S\to\bigl(M(c,b)\to\neg\chi(b)\bigr).
\]
The equality atom occurring inside \(\varphi\) is never substituted: it was removed from the quantified consequent by the propositional weakening step.

Under \(\varphi\), the point condition \(P\) gives \(M(c,a)\). Direct rewriting under \(E(a,b)\) gives \(M(c,b)\), and the displayed consequence of \(S\) contradicts \(\chi(b)\). Therefore
\[
 \varphi\to(S\to\neg P)
\]
is derivable. Generalize at \(c\), use universal distribution, and apply the available diagonal axiom
\(\Gamma(\varphi)=\varphi\to U(c,\varphi)\). This gives
\[
 \varphi\to U(c,S\to\neg P)=\varphi\to Q.
\]
Pointed separation gives \(\neg Q\). Hence \(\neg\varphi\), and classical double-negation elimination gives \(C\). This proof avoids a circular appeal to equality vacuity for \(C\).

For \(a=b\), the desired reverse rewrite is propositional. The cases \(d=a\) reduce, as in the original proof, to direct rewrites and the already established diagonal case \(\chi(t)=M(t,t)\). Thus for all labels the two unrestricted membership rewrites in both directions are recovered:
\begin{equation}\label{eq:R-interface}
 \begin{split}
 E(a,b)&\to\bigl(M(d,a)\leftrightarrow M(d,b)\bigr),\\
 E(a,b)&\to\bigl(M(a,d)\leftrightarrow M(b,d)\bigr).
 \end{split}
\end{equation}

\subsection{Recovering the neighborhood-equality interface}
Define
\[
 R(a,b,z)=U\bigl(z,M(a,z)\leftrightarrow M(b,z)\bigr)
 \wedge U\bigl(z,M(z,a)\leftrightarrow M(z,b)\bigr).
\]
Let \(z_0=\ulcorner E(a,b)\urcorner\). Then \(z_0>a,b\), and the diagonal clause for \(E(a,b)\) gives
\[
 E(a,b)\to U(z_0,E(a,b)).
\]
Lifting~\eqref{eq:R-interface} at \(z_0\) gives
\(E(a,b)\to R(a,b,z_0)\).

It remains to recover the original interface at \emph{every} \(z\ne a,b\), not just at a numerically large binder. If \(z=z_0\), nothing further is required. Otherwise consider either quantifier-free membership matrix
\[
 P(t)=M(a,t)\leftrightarrow M(b,t)
 \quad\text{or}\quad
 P(t)=M(t,a)\leftrightarrow M(t,b).
\]
The two-direction rewrites imply
\(E(z,z_0)\to(P(z_0)\to P(z))\). The witness
\(\exists x_{z_0}E(z,z_0)\), universal distribution, and membership-only vacuity give
\[
 U(z_0,P(z_0))\to P(z).
\]
The antecedent is itself membership-only and vacuous at \(z\), so generalization and distribution yield
\[
 U(z_0,P(z_0))\to U(z,P(z)).
\]
Apply this to both matrices. Consequently
\[
 E(a,b)\to R(a,b,z)\qquad(z\ne a,b)
\]
is derivable without unrestricted vacuity for primitive equality.

The equality-witness head is unchanged. Existential monotonicity and the recovered neighborhood implication supply its interpreted witness form. These are exactly the original completeness proof's three interfaces: two-direction membership rewriting, neighborhood equality, and an equality witness, conventionally denoted (R), (N), and (W).

\subsection{PA calibration and logical scope}
The original PA-refutation compiler uses primitive equality through those interfaces, together with the old gate and the unchanged heads. In the inspected logical and arithmetic-certificate source closure, the uses of equality vacuity occur in the reverse-rewriting and neighborhood modules. Substitute the finite derivations above for those affected templates. Code formation, choosing \(c=\ulcorner\neg C\urcorner\), and \(z_0=\ulcorner E(a,b)\urcorner\), the label-equality case splits, and the finite proof concatenations are primitive recursive. Their correctness is provable in PA. This gives the internal refutation-to-halting direction for the revised calculus.

In the other direction, each new diagonal clause has a uniformly generated first-order proof under the arithmetic interpretation, by the code-freshness lemma. Pointed separation is still supplied by full PA induction as in the inherited arithmetic-calibration argument. The annotated proof-store argument and positive-length tape simulation then give the internal halting-to-refutation direction.

Thus the written calibration argument establishes the same target:
\[
 \PA\vdash\bigl(\neg H_{\Gamma}\leftrightarrow\Con(\PA)\bigr)
\]
for a correct literal implementation of the revised program. This is a calibration to PA; merely replacing the searched system by iPRA would not supply it. The construction applies the general strategy of replacing a primitive schema by smaller derivably equivalent interfaces. Its decisive freshness lemma comes from PA196's existing numeric coding: the magnitude of a formula code guarantees a fresh variable label. We do not attribute that specific diagonal trick to the ConciseIPRA proposals.

